% ===================================================================
%  BAGAN No. 13 — Hotel. --- Rumah makan. Kedai kopi.
%
%  Traduction, faite en 2026, en indonesien d'aujourd'hui, sur l'ido
%  de texto/io. Regles de la colonne : voir 10-bagan-01.tex. Une seule
%  numerotation. Deux notes d'auteur, toutes deux en gras-italique
%  dans le fac-simile ; le bloc t13-08-1 se poursuit apres la seconde
%  sous la MEME cle.
%
%  catur (78) EST UN MOT SANSKRIT ARRIVE PAR LA MER, ET IL RENCONTRE
%  ICI LE MEME MOT VENU PAR LA TERRE. Le releve avait note dans la
%  colonne tamoule que சதுரங்கம், « les quatre corps d'armee », est
%  l'ancetre du shatranj persan et du chess anglais. L'indonesien dit
%  catur, qui est le meme chaturanga, mais arrive avec l'hindouisme au
%  premier millenaire — pas par l'Iran et l'Europe. Sur un salon de
%  cafe ou tout le reste vient d'Occident, un seul objet fait le
%  chemin inverse, et il le fait pour la deuxieme fois dans le releve
%  par une route entierement differente.
%
%  ET LES DEUX AUTRES JEUX DISENT L'HISTOIRE DU BATIMENT OU ILS SE
%  JOUENT. dam les dames (79) vient du neerlandais dammen ; kartu les
%  cartes (82) du portugais carta ; triktrak et biliar sont des
%  emprunts recents. Quatre jeux sur une meme table, quatre couches
%  differentes de la langue — sanskrite, neerlandaise, portugaise,
%  moderne. Le tableau 11 avait etabli qu'il y en a cinq ; celui-ci en
%  montre quatre dans un seul alinea.
%
%  LE PERSONNEL EST MALAIS, COMME EN OURDOU ET POUR LA MEME RAISON.
%  Onze metiers de service, onze mots de fond : pelayan kamar la femme
%  de chambre, kasir la caissiere, penjaga pintu le portier, pesuruh
%  le chasseur, kuli le porteur, tukang cuci la blanchisseuse, kepala
%  pelayan le maitre d'hotel, juru masak le cuisinier, pembantu dapur
%  le marmiton, pelayan le garcon, pengurus le gerant. La colonne
%  ourdoue avait le meme resultat, et la raison est la meme : ces
%  metiers-la s'exercaient dans la langue du pays bien avant qu'un
%  hotel europeen n'ouvre.
%
%  pesuruh, LE CHASSEUR (19 et 54), MERITE D'ETRE COMPARE A چھوٹا.
%  L'ourdou nomme le gamin en livree par son AGE — « le petit ».
%  L'indonesien le nomme par ce qu'on lui fait faire — « celui qu'on
%  envoie ». Deux langues, deux facons d'eviter un nom de fonction que
%  ni l'une ni l'autre n'a, et deux portraits differents du meme
%  enfant.
% ===================================================================

%%K t13-noto-1 noto
\VUnotes{143.2mm}{%
\textit{\VUgras{(*) Untuk perincian tentang kamar, lihat bagan No.
6.}}%
}

%%K t13-apar-1 apar
\VUsaut{3.0mm}
\VUtitre{15.0pt}{60}{BAGAN No. 13}
\VUsaut{1.6mm}
\VUpk{10.2pt}{40}{Hotel. --- Rumah makan.}
\VUsaut{2.0mm}
\VUpk{10.2pt}{40}{Kedai kopi.}
\VUsaut{3.4mm}

%%K t13-01-1 p
1. --- Setelah tiba kemarin sore di Royan, saya menginap di «
\textit{Hotel Samudra} », yang direkomendasikan kepada saya oleh
seorang teman.

%%K t13-01-2 p
Saya diberi sebuah \VUgras{kamar} \textsuperscript{(2)} yang indah, di
\VUgras{lantai} \textsuperscript{(1)} dua, tempat saya tidur dengan
sangat nyenyak (*).

%%K t13-02-1 p
2. --- Pagi ini, keluar dari kamar saya, tempat \VUgras{pelayan kamar}
\textsuperscript{(3)} telah membawa sarapan, saya masuk ke \VUgras{ruang
merokok} \textsuperscript{(5)} yang juga dipakai sebagai \VUgras{ruang
baca} \textsuperscript{(4)}, untuk membaca \VUgras{koran}
\textsuperscript{(6)} sambil mengisap \VUgras{rokok}
\textsuperscript{(8)}. Sesudah membaca, saya turun dengan \VUgras{lift}
\textsuperscript{(9)} dan pergi berjalan-jalan di kota.

%%K t13-03-1 p
3. --- Karena saya lupa menanyakan kepada \VUgras{juru tulis}
\textsuperscript{(12)} hotel pukul berapa makanan dihidangkan di
\VUgras{meja bersama} \textsuperscript{(11)}, saya kembali lebih awal
dan memanfaatkannya untuk pergi mandi di \VUgras{kamar mandi}
\textsuperscript{(10)} hotel. Kemudian saya pergi lagi ke
\VUgras{kasir} \textsuperscript{(15)} untuk menelepon salah seorang
teman saya. Tetapi \VUgras{telepon} \textsuperscript{(13)} sedang
dipakai ; melihat kebingungan saya, \VUgras{kasir perempuan}
\textsuperscript{(14)}, yang sedang mencatat sebuah \VUgras{tagihan}
\textsuperscript{(17)} di \VUgras{buku tamu} \textsuperscript{(16)},
meminta \VUgras{penjaga pintu} \textsuperscript{(18)} menyuruh
\VUgras{pesuruh} \textsuperscript{(19)} mengerjakan urusan saya
\VUgras{dengan sepeda} \textsuperscript{(20)}.

%%K t13-04-1 p
4. --- Saya berterima kasih kepadanya dan keluar untuk memberi persen
kepada \VUgras{kuli} \textsuperscript{(24)} (buruh angkut) yang telah
membawa dari stasiun dengan \VUgras{gerobak tangannya}
\textsuperscript{(25)} \VUgras{kotak topi} \textsuperscript{(26)} saya,
\VUgras{koper saya} \textsuperscript{(27)} dan \VUgras{tas perjalanan
saya} \textsuperscript{(28)}. \VUgras{Pengantar surat}
\textsuperscript{(21)}, yang baru saja tiba, menyerahkan kepada saya
sepucuk \VUgras{surat} \textsuperscript{(22)}.

%%K t13-05-1 p
5. --- Saya masih tinggal beberapa saat di ambang pintu, di dekat
\VUgras{kotak surat} \textsuperscript{(23)}, untuk memandang lalu lalang
di jalan. \VUgras{Kusir} \textsuperscript{(30)} \VUgras{omnibus}
\textsuperscript{(29)} sedang memuat ke atas keretanya \VUgras{koper
besar} \textsuperscript{(31)} seorang \VUgras{penumpang}
\textsuperscript{(32)} yang diantar \VUgras{pemilik hotel}
\textsuperscript{(33)}. Seorang \VUgras{petugas telegraf}
\textsuperscript{(35)} muda dengan tenang membawa sebuah
\VUgras{telegram} \textsuperscript{(34)} untuk seorang tamu hotel ;
seorang \VUgras{wisatawan} \textsuperscript{(36)} dengan
\VUgras{kameranya} \textsuperscript{(38)} tersandang di bahu sedang
membuka \VUgras{buku panduannya} \textsuperscript{(37)}.

%%K t13-tit-1 sub
\VUsaut{4.0mm}
\VUpk{10.2pt}{40}{Waktu makan siang.}
\VUsaut{3.0mm}

%%K t13-06-1 p
6. --- Di depan pagar, seorang \VUgras{pesuruh jalanan}
\textsuperscript{(39)} yang tidak peduli mengantuk di antara
\VUgras{kait pemikulnya} \textsuperscript{(40)} dan \VUgras{kotak
semirnya} \textsuperscript{(41)}, sementara seorang \VUgras{tukang cuci
perempuan} \textsuperscript{(42)} muda, yang membawa sebuah
\VUgras{keranjang} \textsuperscript{(43)} berisi kain, tertawa melihat
wajah sungguh-sungguh \VUgras{kepala pelayan} \textsuperscript{(44)}
yang berdiri di dekat pintu.

%%K t13-07-1 p
7. --- Melihat \VUgras{juru masak} \textsuperscript{(45)}, yang keluar
dari \VUgras{dapurnya} \textsuperscript{(47)} di \VUgras{lantai bawah
tanah} \textsuperscript{(48)} untuk menyuruh \VUgras{pembantu dapur}
\textsuperscript{(46)} melakukan sesuatu, saya teringat bahwa waktu
makan siang sudah dekat, dan saya masuk ke rumah makan, melewati
halaman tempat seorang \VUgras{sopir} \textsuperscript{(50)}
memarkir \VUgras{mobilnya} \textsuperscript{(49)}, sementara seorang
\VUgras{montir} \textsuperscript{(52)} memperbaiki, menurut petunjuk
\VUgras{pengendara sepeda} \textsuperscript{(53)}, sebuah \VUgras{roda
tiga bermotor} \textsuperscript{(51)} yang baru saja mengalami
kecelakaan.

%%K t13-08-1 p
8. --- Saya bertanya kepada seorang \VUgras{pesuruh}
\textsuperscript{(54)} muda yang membawa \VUgras{gelas-gelas bir}
\textsuperscript{(55)} apakah meja bersama ada di bawah
\VUgras{beranda,} dan karena jawabannya membenarkan, saya mengambil
tempat di salah satu meja dan minta dihidangkan makanan yang sangat
baik.

%%K t13-noto-2 noto
\VUnotes{143.2mm}{%
\textit{\VUgras{(*) Dengan bantuan daftar hidangan yang dimuat dalam
kamus, guru akan dapat dengan mudah mengubah-ubah daftar makanan di
bawah ini.}}%
}

%%K t13-tit-2 sub
\VUsaut{4.0mm}
\VUpk{10.2pt}{40}{Makan siang yang lezat.}
\VUsaut{3.0mm}

%%K t13-09-1 p
9. --- Pelayan membawakan kepada saya daftar makanan (*) untuk makan
siang dan daftar anggur. Saya memilih dan memesan sebagai
\VUgras{pembuka} \VUgras{udang} dengan \VUgras{mentega,} dan sebagai
\VUgras{hidangan pertama,} \VUgras{babat cara Caen.} Saya tidak makan
\VUgras{ikan,} melainkan meminta sepotong \VUgras{paha} anak domba, dan
sebagai \VUgras{hidangan sayuran,} \VUgras{bayam} dengan krim.
Kemudian, karena tidak satu pun \VUgras{hidangan antara} yang menarik
hati saya, saya minta dibawakan bermacam-macam hidangan penutup :
\VUgras{keju,} \VUgras{buah} dan \VUgras{biskuit.} Saya menambahkan
lagi \VUgras{anggur} Saint-Emilion yang sangat baik, dan dengan sangat
puas atas santapan saya, saya meninggalkan meja sesudah memberi kepada
\VUgras{pelayan} \textsuperscript{(57)} \VUgras{persen}
\textsuperscript{(59)} yang lumayan.

%%K t13-10-1 p
10. --- Melihat saya siap pergi, \VUgras{pengurus}
\textsuperscript{(60)} mengusulkan agar saya naik ke balkon untuk
mengagumi pemandangan \VUgras{Samudra} \textsuperscript{(62)} yang
megah. Di kejauhan mulai tampak \VUgras{perahu-perahu nelayan kecil}
\textsuperscript{(63)}, \VUgras{kapal-kapal layar}
\textsuperscript{(64)} dan sebuah \VUgras{kapal pos}
\textsuperscript{(65)} yang sangat besar, yang bayangan hitamnya
menonjol pada \VUgras{awan-awan} \textsuperscript{(67)} putih di
\VUgras{kaki langit} \textsuperscript{(66)}. Angin ringan membuat
berkibar \VUgras{panji} \textsuperscript{(70)} yang ditancapkan di atas
\VUgras{lambang} \textsuperscript{(69)} \VUgras{aula}
\textsuperscript{(68)}.

%%K t13-tit-3 sub
\VUsaut{3.0mm}
\VUpk{10.2pt}{40}{Hiburan.}
\VUsaut{3.0mm}

%%K t13-11-1 p
11. --- Pemandangan itu begitu menarik sehingga saya memutuskan besok
akan naik ke teras kedai kopi sambil membawa \VUgras{teropong ganda}
\textsuperscript{(71)} atau \VUgras{teropong} \textsuperscript{(72)}.

%%K t13-12-1 p
12. --- Meninggalkan balkon, saya pergi ke kedai kopi hotel untuk
menikmati sebuah \VUgras{minuman} \textsuperscript{(73)} sambil bermain
\VUgras{biliar} \textsuperscript{(75)}. Tanpa berhenti saya melintasi
ruang kedai kopi tempat banyak tamu bermain \VUgras{catur}
\textsuperscript{(78)}, \VUgras{dam} \textsuperscript{(79)},
\VUgras{domino} \textsuperscript{(80)}, \VUgras{triktrak}
\textsuperscript{(81)} atau \VUgras{kartu} \textsuperscript{(82)}, dan
saya langsung naik ke \VUgras{ruang biliar} \textsuperscript{(74)}.

%%K t13-13-1 p
13. --- Ketika permainan selesai, saya turun ke lantai dasar dan
meminta kepada pelayan sebuah \VUgras{amplop} \textsuperscript{(84)},
\VUgras{kertas} \textsuperscript{(83)}, sebuah \VUgras{perangko}
\textsuperscript{(85)}, sebuah \VUgras{pena} \textsuperscript{(87)} dan
\VUgras{tinta} \textsuperscript{(86)} untuk menulis surat.

%%K t13-13-2 p
Kemudian saya memanggil seorang \VUgras{kusir} \textsuperscript{(92)}
yang \VUgras{keretanya} \textsuperscript{(88)} berhenti di depan kedai
kopi. Saya menanyakan harganya menurut jam atau sekali jalan, dan
ketika kami sepakat tentang harganya, ia membukakan untuk saya
\VUgras{pintu kereta} \textsuperscript{(89)} dan mempersilakan saya
duduk. Ia naik kembali ke \VUgras{tempat duduknya}
\textsuperscript{(91)}, dengan \VUgras{cambuk} \textsuperscript{(93)}
menjalankan kuda-kuda dengan cepat, dan kami berangkat untuk sebuah
tamasya yang memesona.

\VUsaut{6.0mm}
\VUfilet{22mm}
